%
%
% report_2025_1_intro.tex
%
%
%
\documentclass[jctc-draft.tex]{subfiles}
\begin{document}
\section{Introduction}
Since the early conception of quantum computers, quantum chemistry simulation
has been recognized as a highly promising application area expected to be
significantly advanced by quantum computing
\cite{Cao2019.doi:10.1021/acs.chemrev.8b00803}. Currently, mainstream quantum
hardware is inherently prone to computational errors due to noise, and the
practical implementation of error correction technologies is ongoing. The
surface code method, a promising error correction technique for quantum
computers, requires a significantly higher level of redundancy compared to
classical error correction methods \cite{Fowler2012.PhysRevA.86.032324}.
Consequently, realizing a practical fault-tolerant quantum computer (FTQC) faces
implementation challenges, primarily the need to vastly increase the number of
qubits on hardware devices, and is expected to take several more years for
development. In the current Noisy Intermediate-Scale Quantum (NISQ) era, where
error correction is not yet implemented and qubit counts are limited,
Variational Quantum Eigensolver (VQE)-based algorithms utilizing second
quantization are the mainstream of quantum chemistry research
\cite{peruzzo2014variational}. VQE combines a quantum computer with a classical
computer to reduce the required number of qubits and circuit depth, making it
executable on NISQ devices as it can tolerate the presence of noise. However,
once FTQC is realized---overcoming the impact of noise--- and qubit count
constraints being relaxed, a broader range of algorithms is expected to be
researched and put into practical use.

The grid-based approach, a representative method for quantum chemistry
simulation based on first quantization, is one of the ultimate targets in this
field because it can directly handle electron interactions and also handle
electron structure dynamics because nuclei-centered orbital basis functions are
not used. However, it requires substantially more computational resources than
VQE \cite{kassal2008.pnas.0808245105}. In the grid method, the indices of
equally spaced grid points are mapped to the tensor product of computational
basis states of multi-qubit registers each corresponding to spatial coordinates.
The values of the wavefunction at these grid points are then stored as the
amplitudes of the corresponding basis to simulate its time evolution. While
algorithms based on the Suzuki-Trotter decomposition are widely known for time
evolution calculations, algorithms employing more sophisticated approaches have
also been proposed \cite{Berry2015.PhysRevLett.114.090502, Babbush_2018,
YuanSu.PRXQuantum.2.040332, eklund2026.arXiv2603.19007}. Nevertheless, because
all of these grid-based algorithms generally require a large number of qubits
and significant circuit depth, and cannot tolerate computational errors, reports
of their operational verification are strictly limited. To the best of the
authors' knowledge, no calculations using actual quantum computer hardware have
been reported to date; existing reports have relied on quantum computer emulators
running on classical machines. For example, Benenti et al. performed
calculations using a simplified potential function
\cite{Benenti2008.10.1119/1.2894532}. Chan et al. utilized both a quantum
computer emulator and a classical program, processing the Coulomb potential
calculations on the classical side \cite{Chan2023.sciadv.abo7484}. In other
words, an example of computing the Coulomb potential entirely within a quantum
circuit has not yet been reported, even on an emulator. Therefore, in this
study, we present the configuration and operational demonstration of a
grid-based simulator implemented entirely with quantum circuits, including the
Coulomb potential calculation. The software used to generate the simulator
circuits is crsQ, developed by the author
\cite{takahashi2024.doi:10.1021/acs.jctc.4c00708}. We implement the wave
function time evolution calculation via the grid method entirely as a quantum
circuit, execute it on a quantum computer emulator, and evaluate the
computational results. To achieve this, we address the following three
challenges.

The first challenge is to present a circuit configuration that can operate on an
emulator with a limited number of qubits. In this study, we simplify the
computational model to fit within the emulator's constraints, targeting 1D and
2D models of a single hydrogen atom with fixed nuclear coordinates. The
fundamental approach for calculating the Coulomb potential is based on binary
arithmetic operations \cite{Zalka-rspa.1998.0162}. While methods have been
proposed to reduce the gate count by precomputing parts of the calculation on a
classical computer and using a QROM-based lookup table, these approaches require
a significant number of ancilla qubits \cite{Sanders2020.PRXQuantum.1.020312}.
In this study, to conserve qubit resources, we adopt the Uniformly Controlled
Rotation Z (UCRZ) circuit \cite{Mottonen2004.PhysRevLett.93.130502}, which
applies rotations to a target qubit based on the values in a control register.
For time-evolution, we employ the basic Suzuki-Trotter
decomposition \cite{Zalka-rspa.1998.0162}. Through this approach, we provide
circuit configuration examples suitable for emulators and early-stage FTQC
devices. In the grid method, the scale of the wavefunction initialization
circuit also poses a challenge, as preparing an arbitrary state requires a
number of gates that scales exponentially with the grid size. As a means to
reduce this scale, a method has been proposed to convert initial states
expressed as linear combinations of atomic orbitals into a Matrix Product State
(MPS) format---or Tensor-Train format
\cite{Oseledets2011.doi:10.1137/090752286}---and construct the initialization
circuit based on low-rank approximations \cite{Huggins2025.PRXQuantum.6.020319}.
This technique can significantly reduce the circuit scale by utilizing a bond
rank that corresponds to the smoothness of the wavefunction and the strength of
inter-particle correlations. However, since this study focuses on small-scale
grid sizes feasible for emulation, we utilize simple initialization circuits
without MPS compression to prioritize ease of implementation. 

The second challenge is to determine the handling of the singularity (pole) of
the Coulomb potential. In the grid method, it is necessary to calculate the
potential energy for all possible combinations of coordinates for all particles
and update the phase of the wavefunction, that is proportional to the energy.
The point where the coordinates of the nucleus and the electron coincide is a
singularity of the Coulomb potential, where the value diverges if left
unaddressed. It is known that methods using a plane-wave basis (momentum-space
representation) can avoid direct calculations of the singularity in real space
by treating the potential as an operation in momentum space \cite{Babbush_2019,
eklund2026.arXiv2603.19007}. However, these approaches require a significant
number of ancilla qubits, making them unsuitable for execution on small-scale
quantum computers. In contrast, while the position-basis approach can be
implemented with fewer ancilla qubits, it necessitates an explicit strategy for
addressing the singularity. To this end, several techniques have been proposed,
such as employing a soft-core potential that adds a fixed offset to the
inter-particle distance \cite{Kivlichan_2017}, or introducing a relative shift
between the grids for the nucleus and electrons to prevent the distance from
becoming zero \cite{Chan2023.sciadv.abo7484}. In this study, we propose a method
to determine an effective potential value based on the integral average (or
volume average) of the potential function within the grid cell containing the
singularity, and we verify its validity through numerical evaluation.

The third challenge is to establish a method for determining simulation
parameters that minimize error within a given number of qubits. In the grid
method, it is necessary to first define simulation parameters such as the
spatial resolution (grid spacing) and the size of the simulation domain.
Furthermore, when using the Suzuki-Trotter decomposition, the time step size
must also be determined. The Suzuki-Trotter decomposition approximates the
exponential of a Hamiltonian---which contains inherently non-commuting potential
and kinetic energy terms---as a product of the individual exponentials. This
decomposition introduces errors arising from the non-commutativity of the
operators. Early error analyses were based on operator norms, evaluating the
worst-case scenario across all possible input states. As a result, these
estimates were noted to be several orders of magnitude larger than the errors
empirically observed in practical numerical calculations. However, by focusing
on commutator scaling, more tight error bounds have become possible
\cite{Childs2021.PhysRevX.11.011020}. Subsequently, it was shown that the bound
can be further tightened by considering state-dependency rather than relying on
worst-case evaluations \cite{Burgarth2024.PhysRevResearch.6.043155}. That study
demonstrated that when the wavefunction value is non-zero at the potential
singularity (e.g., for $s$-orbitals), the error exhibits a sublinear scaling of
$n^{-1/4}$ with respect to the number of Trotter steps $n$. It also showed that
higher-order Trotter formulas do not contribute to error reduction for such
states. This result was extended to multi-particle systems, proving that for
unbounded Hamiltonians like the Coulomb potential, the dependence on the number
of steps $n$ is $n^{-1/4}$, and the error increase relative to the particle
number is polynomial \cite{fang2026trotter}. In our study, through numerical
evaluations based on an emulator and classical programs, we present an
integrated approach for a single-electron model: we optimize spatial parameters
by balancing the tradeoff between spatial discretization error and domain
truncation error, and we determine time parameters by considering the
requirements of the sampling theorem alongside recent Trotter error evaluations.

\end{document}